\chapter{Analisi delle tecnologie e metodologie impiegate}

In questo capitolo verranno fornite informazioni inerenti alle principali tecnologie e metodologie utilizzate durante lo sviluppo di Sapienza-DC, in modo tale da chiarire il ruolo e il funzionamento di ciascuno strumento impiegato durante la fase di implementazione del progetto.

\section{Concetti preliminari di NLP}
L'NLP (Natural Language Processing) è una branca dell'Intelligenza Artificiale che studia la lettura, comprensione, interpretazione nonché generazione, da parte di un calcolatore, di linguaggio umano, in modo naturale.

NLP combina in un unicum la linguistica computazionale, ossia le regole del linguaggio, con tecniche avanzate di machine learning\footnote{machine learning: branca dell'IA avente come obiettivo quello di apprendere autonomamente da un insieme di dati in modo tale da essere in grado successivamente di effettuare previsioni su nuovi dati. } e deep learning\footnote{deep learning: sottoinsieme del machine learning in cui si utilizzano reti neurali multistrato costituite da neuroni artificiali interconnessi. }. In una tipica pipeline NLP, il testo viene dapprima elaborato mediante tokenizzazione, la quale consiste nella suddivisione del testo in unità più piccole, quali parole o frasi. Dopodiché, ciascun token viene standardizzato, mediante conversione in minuscolo di ciascun carattere, la rimozione delle cosiddette stop word (parole del linguaggio comune utilizzate frequentemente, come ``è'', ``la'', etc.) nonché caratteri speciali o di punteggiatura, considerati come rumore in questa fase iniziale.

Successivamente, ciascun token viene trasformato, mediante opportune tecniche di NLP, in rappresentazioni numeriche su cui un calcolatore può operare. Ad esempio, tramite opportuni modelli di embedding (incorporamento) è possibile trasformare un token in un vettore n-dimensionale di reali, ove è possibile memorizzare il significato semantico del token originale. Questo, peraltro, sarà proprio l'approccio utilizzato da Sapienza-DC, come vedremo più in avanti.

A partire dai dati così ottenuti, vengono estratte ulteriori informazioni mediante opportune tecniche computazionali, al fine di identificare ruoli grammaticali delle parole, argomenti e sentiment, vale a dire il tono emotivo del testo. Infine, le nozioni raccolte vengono utilizzate per addestrare modelli di machine learning, capaci di apprendere tali relazioni e utilizzarle una volta completato l'addestramento per generare linguaggio naturale.         

\subsection{Large Language Model (LLM)}
Un LLM è un particolare tipo di modello linguistico capace di generare linguaggio naturale mediante addestramento su una grande quantità di dati. Il termine ``large'' deriva dal grande (spesso nell'ordine di miliardi) numero di parametri del modello probabilistico appresi in fase di training. Gli LLM sono basati su deep learning e utilizzano come componenti fondamentali reti neurali e trasformatori\footnote {trasformatore: particolare architettura di rete neurale artificiale che permette l'elaborazione parallela di diverse porzioni di testo, cogliendo inoltre eventuali relazioni tra parole distanti nel testo grazie a un meccanismo noto come self-attention.} (transformers).

Un LLM genera testo in modo puramente statistico, infatti, se \((w_1, w_2,\, \dots \,, w_T)\) sono parole che costituiscono una frase $f$, allora la probabilità che l'LLM $\mathscr{L}$ generi $f$ è data da: 
\[\mathbb{P}(w_1, w_2,\, \dots \,, w_T) = \prod_{j=1}^{T} \mathbb{P}(w_j \, | \, w_1, w_2,\, \dots \,, w_{j-1})\]
dove la probabilità di generare la j-esima parola $w_j$ dipende dai parametri appresi durante l'addestramento del modello. La generazione, detta autoregressiva, avviene un token alla volta ed è pertanto sequenziale:  ciascun nuovo token dipenderà direttamente da quelli generati in precedenza. Vi sarebbero numerosissime tematiche da approfondire a tal riguardo, ma, per brevità, considereremo sufficienti queste nozioni intuitive.  

\subsection{Retrieval-Augmented Generation (RAG)}
Retrieval-Augmented Generation, ossia la generazione potenziata mediante recupero di dati, è un approccio che consiste nell'ottimizzazione dell'output di un LLM.

RAG prevede l'estrazione di informazioni e conoscenze utili da fonti autorevoli, al fine di accrescere le nozioni possedute dal modello. Tale strategia diviene estremamente utile quando è necessario che l'LLM fornisca una risposta a query altamente specifiche o inerenti a fatti recenti, altrimenti non contenuti nel dataset di addestramento dell'LLM.

Questa tecnica è adottata da Sapienza-DC per estrarre contenuti ed informazioni domain-specific con lo scopo di generare una risposta alle query proposte al chatbot da parte dell'utente. 

\begin{figure}[htbp]
    \centering
    \includegraphics[width=0.9\textwidth]{rag_diagram}
    \caption[Elementi fondamentali di uno schema di RAG.]{Elementi fondamentali di uno schema di RAG.\footnotemark}
    \label{fig:diagramma_rag}
\end{figure}
\footnotetext{Fonte immagine: \href{https://obot.ai/wp-content/uploads/2025/09/RAG_1200_x_800px_c15f070854.png}{obot.ai}. Rilasciata sotto la seguente licenza: \href{https://creativecommons.org/licenses/by/4.0/deed.it}{CC BY 4.0}.}

\subsection{Bi-encoder, cross-encoder}
Un bi-encoder (o dual-encoder) è un modello utilizzato per convertire testo in vettori di reali, detti embedding dell'informazione di partenza. Il nome deriva dall'utilizzo di due reti neurali separate aventi lo scopo di trasformare rispettivamente query e documenti di riferimento nei relativi embedding.

L'idea a questo punto è molto semplice: verificare, mediante opportune operazioni sui vettori così ottenuti, quali documenti, o porzioni di essi, risultino più semanticamente affini alla query fornita dallo user. Un approccio utile a questo scopo che vedremo successivamente consisterà nel calcolo della cosine similarity tra $v_q$ e $v_c$, vettori rispettivamente della query $q$ e del chunk $c$, appartenente a un determinato documento ottenuto dalla corrispettiva pagina web.

I bi-encoder sono notoriamente veloci, tuttavia, analizzando separatamente query e chunk dei documenti non godono di un'accuratezza eccellente. Per questo motivo è necessario utilizzare anche i cross-encoder. Questi ultimi sono modelli aventi lo stesso identico scopo, ma che offrono una precisione molto maggiore dei dual-encoder, a discapito di un maggiore costo computazionale.

In particolare, i cross-encoder devono necessariamente elaborare contemporaneamente query e chunk di testo, sfruttando i meccanismi di attenzione offerti dai transformer affinché sia colto il contesto della domanda e la reale pertinenza tra quest'ultima e la porzione di documento sotto esame. L'output di un cross-encoder non consiste in un vettore, ma in uno score, indice del grado di rilevanza del chunk $c$ rispetto alla query $q$.

Nella fase di implementazione, utilizzeremo i cross-encoder, in seguito a un filtraggio preliminare mediante bi-encoder e cosine similarity, per effettuare reranking (riordinamento) dei chunk vari chunk estratti dalle fonti di interesse.

\section{Framework, strumenti e librerie di sviluppo}

\subsection{FastAPI}
FastAPI~\cite{fastapi_docs} è un framework web moderno, open source e ad alte prestazioni, progettato per la sviluppo di API\footnote{API - Application Programming Interface: un insieme di regole e protocolli che permettono a due applicazioni di comunicare correttamente.} (Application Programming Interface) REST (REpresentational State Transfer) in Python.

Rilasciato nel 2018, FastAPI è ben presto divenuto uno degli standard per lo sviluppo backend in Python, grazie alla sua efficienza nonché facilità di utilizzo.

Il framework usufruisce di due librerie fondamentali quali Starlette, un toolkit web ASGI (Asynchronous Server Gateway Interface) per la gestione di servizi asincroni, e Pydantic, grazie a cui è possibile una validazione automatica e personalizzabile dei dati e la creazione di modelli Pydantic per la definizione dei tipi di input e output di ciascun endpoint\footnote{endpoint: nell'ambito delle API, un URL specifico tramite cui un'applicazione riceve ed elabora richieste API}.

Tali funzionalità facilitano e rendono estremamente intuitivo lo sviluppo di backend, garantendo prestazioni equiparabili ad altre tecnologie popolari come Node.js per lo sviluppo server-side.

Tra le funzionalità più apprezzate di FastAPI risalta la generazione automatica della documentazione API. Ogni qualvolta viene definito un nuovo endpoint, FastAPI genera automaticamente documentazione e interfacce Web UI (Swagger UI) mediante cui inviare richieste ai fini, ad esempio, di testing dell'applicazione.

Infine, il framework segue e adotta i principali standard web moderni per lo sviluppo di API basate su HTTP (\textit{Hypertext Transfer Protocol}), quali OpenAPI e JSON Schema per la validazione nonché serializzazione dei dati scambiati.    

\subsection{LangChain}
LangChain~\cite{langchain_docs} è un framework open source progettato per facilitare lo sviluppo di applicazioni che fanno uso di LLM.

Gli impieghi possibili di LangChain sono numerosi, ma ai fini dello sviluppo di Sapienza-DC ci limiteremo a utilizzare la classe RecursiveCharacterTextSplitter per l'estrazione di chunk coerenti dalle fonti di riferimento, preservando la semantica delle informazioni originali.   

\subsection{Hugging Face}
Hugging Face~\cite{huggingface_docs} è una piattaforma open source dedicata all'intelligenza artificiale e al machine learning. Grazie alla libreria Model Hub è possibile il download di modelli già addestrati per svariati scopi, tra cui quelli di cui usufruiremo ai fini dello sviluppo del progetto.

\subsubsection{FastEmbed}
FastEmbed~\cite{fastembed_docs} è una libreria Python adatta a generare embedding di testi e immagini. Come già accennato in precedenza l'embedding permette di ottenere un vettore di reali che mantengano la semantica linguistica, nel caso di testo, del dato originale.

Tra gli innumerevoli modelli offerti da FastEmbed si è scelto di utilizzare intfloat/multilingual-e5-small come bi-encoder, il cui download è possibile grazie alla Model Hub di Hugging Face.

\subsubsection{SentenceTransformers}
SentenceTransformers~\cite{sentencetransformers_docs} è un framework ideato per la computazione di embedding e l'utilizzo di modelli di reranking. Utilizzeremo quest'ultimo per poter disporre di ms-marco-MiniLM-L-6-v2, un cross-encoder adatto ai nostri scopi. Anche in questo caso il download del modello avverrà automaticamente tramite Hugging Face durante l'inizializzazione della nostra applicazione.


\subsection{Ollama e llama.cpp}
Ollama~\cite{ollama_docs} è un framework open source concepito per semplificare l'esecuzione, la configurazione e la gestione locale degli LLM su sistemi operativi come macOS, Windows e Linux.

Il framework espone una serie di API RESTful standard che consentono di interagire programmaticamente con i modelli scaricati localmente. Grazie al supporto di librerie ufficiali e sviluppate dalla community per linguaggi come Python e JavaScript, è possibile inoltre integrare le capacità di ragionamento degli LLM all'interno di applicazioni, script o flussi di lavoro aziendali in modo estremamente semplice.

Ollama offre, oltretutto, un'immagine Docker ufficiale da eseguire in un opportuno container. Maggiori informazioni su Docker e sul concetto di container verranno fornite nella relativa sezione di questo capitolo.

Si noti per di più che, poiché i modelli sono eseguiti interamente in locale, il sistema non trasmette alcuna informazione a server esterni, garantendo massima privacy e totale indipendenza da connessioni internet o provider di terza parte.

Inoltre, giacché l'elaborazione viene effettuata direttamente sull'hardware dell'utente, il costo di ogni singola query posta all'LLM diviene fondamentalmente nullo.

In aggiunta a ciò, il framework supporta e aggiorna costantemente una libreria di modelli pubblici, tra cui Llama, Mistral, Gemma e Phi-3, permettendo di selezionare l'architettura più idonea alle proprie necessità.

\begin{figure}[htbp]
    \centering
    \includegraphics[width=0.9\textwidth]{diagramma_di_flusso_ollama.png}
    \caption{Diagramma di flusso di un prompt effettuato a un LLM tramite Ollama.}
    \label{fig:diagramma_flusso_ollama}
\end{figure}


Per permettere a reti neurali complesse di funzionare su hardware consumer, Ollama implementa un processo noto come quantizzazione dei modelli. Questa tecnica riduce la precisione matematica dei pesi dell'LLM, abbattendo pertanto in maniera drastica la quantità di RAM nonché la potenza di calcolo richieste per l'esecuzione del modello.

Tuttavia, questo approccio comporta alcuni \textit{trade-off}\footnote{trade-off: compromesso per il quale al fine di ottenere un vantaggio si accetta di cedere o rinunciare a qualcos'altro.}  inevitabili. In primis, a causa dei limiti hardware dei dispositivi consumer, i modelli eseguiti localmente presentano un numero di parametri inferiore e, conseguentemente, una qualità dell'output e capacità di ragionamento più limitate.

In secondo luogo, l'effettiva velocità di esecuzione di Ollama è rigidamente vincolata all'hardware della macchina ospitante, prediligendo architetture basate su GPU a quelle CPU. Naturalmente, l'esecuzione di un LLM tramite Ollama è comunque possibile su macchine sprovviste di GPU, a patto di accettare un'inferenza più lenta del modello.

L'obiettivo di Sapienza-DC consisterà difatti nel permettere a qualsiasi utente, dotato di una scheda video o anche semplicemente del proprio processore, di eseguire il progetto localmente, fornendo la possibilità di affidarsi, in caso di necessità, anche ad API esterne, qualora l'hardware della propria macchina dovesse risultare non sufficientemente performante. Entreremo maggiormente nei dettagli di questa tematica nei capitoli successivi.

A tal proposito, vedremo come sia possibile utilizzare direttamente llama.cpp per gestire manualmente il caricamento, la gestione nonché l'esecuzione di query verso LLM locali, con lo scopo di ridurre l'overhead introdotto da Ollama, il quale non è altro che un wrapper di quest'ultimo.

Brevemente, llama.cpp è un'implementazione in linguaggio C/C++ del codice di inferenza dei modelli LLaMA (Large Language Model Meta AI), successivamente esteso a molteplici altre architetture, il cui scopo fondamentale consiste nel massimizzare la velocità di esecuzione locale di LLM. Rimuovendo l'overhead dovuto a Ollama e accettando la maggiore difficoltà implementativa che consegue da questa scelta, saremo in grado di minimizzare i requisiti hardware richiesti per eseguire modelli localmente, soprattutto su macchine sprovviste di GPU.     

\subsection{SearXNG}
Utilizzato in una prima versione del progetto, SearXNG~\cite{searxng_docs} è un metamotore di ricerca che aggrega risultati da svariati servizi di ricerca e database\footnote{database: base di dati, una raccolta organizzata di informazioni o dati strutturati.}. Garantisce anonimità e privacy per gli utenti che lo utilizzano, i quali non subiscono tracking né profilazione.   

\subsection{Crawl4AI}
Crawl4AI~\cite{crawl4ai_docs} è una libreria Python open source impiegata principalmente per estrarre e ripulire dati in modo strutturato e automatico da pagine web.

In realtà, Sapienza-DC non utilizza direttamente questa libreria, ma fa affidamento a un secondo progetto precedentemente sviluppato, MWP~\cite{mwp_repo} (Minerva Web Parser), per effettuare scraping e parsing di pagine web mediante Crawl4AI. Maggiori informazioni sul funzionamento di MWP verranno presentate nei capitoli successivi.

Minerva Web Parser è stato realizzato come progetto di esonero (al quale si accennava nell'introduzione di questo elaborato) per il corso di Laboratorio di Ingegneria Informatica, ed è stato successivamente aggiornato per includere nuove funzionalità durante lo sviluppo di Sapienza-DC stesso.  
\subsection{Jinja2}
Jinja2~\cite{jinja2_docs} è un engine di templating utilizzato per definire pagine web generiche contenenti opportuni segnaposto da popolare a runtime tramite chiamate API al backend.

Verrà utilizzato per implementare un frontend in grado di servire, previa ottenimento di dati dal server API, pagine HTML (HyperText Markup Language) dinamiche.  
\subsection{TailwindCSS}
TailwindCSS~\cite{tailwindcss_docs} è un framework CSS\footnote{CSS - Cascading Style Sheets: linguaggio utilizzato per definire la formattazione e il layout di documenti HTML.} open source per lo sviluppo di UI moderne ed eleganti. Tailwind mette a disposizione una serie di classi CSS dette di utility da combinare e utilizzare insieme ad altre classi per definire lo stile di ciascun elemento in modo intuitivo e modulare.

\subsubsection{Utility class}

Le utility classes rappresentano l'approccio odierno allo styling CSS: l'idea è quella di avere numerose classi che svolgono compiti ben precisi, come, ad esempio, l'applicazione di una determinata proprietà di stile a un elemento HTML. 

\section{Containerizzazione ed orchestrazione}
La containerizzazione è una tecnologia che consente di racchiudere un'applicazione, assieme alle sue dipendenze, in un'unita standardizzata detta container. I container permettono agli sviluppatori di garantire:
\begin{itemize}
        \item \textbf{portabilità} - affinché l'applicazione possa essere eseguita correttamente a prescindere dall'ambiente o dalla piattaforma su cui viene eseguito il container;
        \item \textbf{isolamento} - ciascun container opera indipendentemente dagli altri, in modo tale da evitare possibili interferenze;
        \item \textbf{efficienza} - i container condividono il kernel del sistema operativo della macchina sulla quale vengono eseguiti e sono pertanto più leggeri delle macchine virtuali, garantendo dunque un minore overhead computazionale.
\end{itemize}
L'orchestrazione di container consiste invece nel processo di automazione della gestione di questi ultimi, al fine di sviluppare applicazioni su larga scala. Un'applicazione può difatti sfruttare molteplici container per implementare diverse funzionalità in modo incredibilmente modulare.

L'impiego di container separati permette inoltre di evitare problematiche relative a SPOF (Single Point Of Failure). Se un servizio non dovesse risultare disponibile sarà il singolo container a non essere raggiungibile, mentre gli altri servizi offerti dall'app rimarranno online, purché non dipendenti dal container momentaneamente fuori uso.   
\subsection{Architettura a microservizi: Docker e Docker Compose}
L'approccio sopra descritto prende il nome di architettura a microservizi. Ribadiamo sinteticamente l'intuizione alla base di quest'ultimo: dividere le funzionalità offerte dalla nostra applicazione in una serie di microservizi ospitati in altrettanti container, opportunamente orchestrati.

Al fine di implementare quanto appena visto nel progetto, faremo uso di Docker~\cite{docker_docs}, una piattaforma open source per la distribuzione automatizzata di applicazioni all'interno di container isolati.

Docker fornisce un ambiente isolato, riproducibile e portabile, pensato per l'esecuzione nonché la distribuzione di un progetto software in un ambiente di testing o di produzione.

Per orchestrare ciascun container viene inoltre reso disponibile Docker Compose, uno strumento per la gestione di applicazioni multi-container basato su file YAML\footnote{YAML - Yet Another Markup Language: un formato di serializzazione di dati basato su linguaggio naturale.}.        
\section{Database relazionale}
Un database relazionale è un particolare tipo di base di dati in cui i dati sono organizzati sotto forma di righe e di colonne. Ciascuna tabella della base di dati rappresenta una determinata relazione $\mathscr{R} \subseteq S_1 \times S_2 \times \dots \times S_n$, cioè un sottoinsieme del prodotto cartesiano di $n$ insiemi. Ciascuna tupla $t \in \mathscr{R}$ sarà nella forma $t = (s_1, s_2, \, \dots \,, s_n)$, dove $s_1 \in S_1, s_2 \in S_2 \, \dots \, , s_n \in S_n$. Nella pratica, ogni riga della tabella consisterà in una $n$-upla di elementi che assumono valori in un dominio (un intero, una stringa. etc.) della corrispettiva colonna, detta attributo o campo. Si noti infine che, essendo $\mathscr{R}$ un insieme, $\neg \exists \,(t,\, \hat{t}\,) \in \mathscr{R} \, | \, t = \hat{t}$, cioè $\mathscr{R}$ non ammette tuple duplicate.
\subsection{SQL}
SQL (Structured Query Language) è il linguaggio utilizzato per interagire programmaticamente con un RDBMS (Relational DataBase Management System), al fine di aggiungere, modificare, eliminare o ricercare informazioni all'interno delle tabelle della base di dati (DML\footnote{DML: Data Manipulation Language}); gestire la struttura del database e delle tabelle che compongono quest'ultimo (DDL\footnote{DDL: Data Definition Language}); la concessione e la revoca di permessi relativi alla base di dati (DCL\footnote{DCL: Data Control Language}); controllo e gestione delle transazioni, ossia insiemi di operazioni che devono essere eseguite in modo atomico (TCL\footnote{TCL: Transaction Control Language}).         
\subsection{MariaDB}
MariaDB~\cite{mariadb_docs} è un sistema di gestione di basi di dati relazionali open source. In particolare, non è altro che un fork di MySQL~\cite{mysql_docs}. Progettato per garantire compatibilità con MySQL stesso, MariaDB offre un'immagine Docker per la gestione l'interazione con database mediante container.

\section{Metodologia di sviluppo agile}
La metodologia agile consiste in un approccio allo sviluppo basato sulla distribuzione di software implementati secondo un processo rapido e iterativo. Contrapponendosi ai modelli tradizionali di sviluppo del software, come ad esempio il modello a cascata (waterfall model\footnote{waterfall model: metodologia di sviluppo software basata sulla definizione precisa e strutturata di un insieme di requisiti iniziali da rispettare, seguita da un'implementazione prevedibile e lineare di questi ultimi.}), la metodologia agile propone un approccio meno strutturato e incentrato sulla consegna ripetuta nel tempo di diversi MVP\footnote{MVP - Minimum Viable Product: una versione basilare di un prodotto o servizio volta ad offrire solo le funzionalità essenziali richieste dal cliente.}.

Questi ultimi vengono sottoposti, volta per volta, ad analisi da parte dei clienti nonché degli sviluppatori con l'obiettivo di rifinire l'approccio utilizzato, modificando anche completamente la soluzione proposta qualora ritenuto necessario, per il raggiungimento di un prodotto migliore nella successiva iterazione. Si tratta di una metodologia frequentemente impiegata in progetti dinamici e adattivi, come, ad esempio, Sapienza-DC.

\section{Paradigma client-server}
Il paradigma client-server rappresenta uno dei modelli più importanti per la comunicazione web e lo sviluppo di sistemi distribuiti. Il modello si fonda sulla separazione netta tra entità che richiedono un certo numero di servizi o risorse, detti client, e i fornitori di questi ultimi, detti server, che si occupano di servire le richieste ricevute dai client.

La comunicazione tra client e server avviene, a livello applicativo, attraverso richieste e risposte HTTP/HTTPS, nei corpi delle quali possono essere inclusi dati, formattati secondo diversi standard, tra cui JSON\footnote{JSON - JavaScript Object Notation: formato testuale di serializzazione dei dati.}, XML\footnote{XML - eXtensible Markup Language: un metalinguaggio per la definizione di linguaggi di markup.} e molti altri. Tale comunicazione avviene solitamente secondo una logica stateless, cioè priva di stato, per la quale il server non mantiene una cronologia delle richieste ricevute precedentemente dal client, ma elabora ognuna di esse in modo indipendente.

Alcuni vantaggi di questo approccio includono, ad esempio, una migliore scalabilità, sicurezza e manutenibilità dei sistemi e delle applicazioni web.

\begin{figure}[H]
    \centering
    \includegraphics[width=1.0\textwidth]{diagramma_architettura_client_server.png}
    \caption{Diagramma\protect\footnotemark\ sull'architettura client-server.}
    \label{fig:diagramma_architettura_client_server}
\end{figure}

\footnotetext{\textbf{N.B.} Immagine di base generata dall'autore mediante IA (\textit{Nano Banana}, modello basato su Google Gemini) in data \textit{4 Giugno 2026}. A tal riguardo, è stato utilizzato il seguente prompt: \textit{«A clean vector illustration of a client-server architecture. On the left, a modern laptop representing the client. On the right, a server rack with a small database icon next to it. In the top center, a simple cloud icon representing the internet. Employ a minimalist tech style, white background, isometric perspective, containing the most relevant facts regarding this type of architecture.»}. L'immagine generata è stata successivamente rielaborata, modificata e annotata manualmente dall'autore.}
